\textbf{Phần 1. Trắc nghiệm lựa chọn}
\setcounterpageref{ex}{0}

\begin{ex} %[3D2V56]
Tính giá trị phép tính: $1425 \times 4$.
\choice
{$5600$}
{$5800$}
{$5680$}
{\True $5700$}
\loigiai{
$1425 \times 4 = 5700$.
\textbf{Đáp án của câu hỏi là D}
}
\end{ex}

\begin{ex} %[3D2V56]
Một xã có 3 thôn, mỗi thôn có 1240 nhân khẩu. Toàn xã đó có bao nhiêu nhân khẩu?
\choice
{$3620$ nhân khẩu}
{\True $3720$ nhân khẩu}
{$3820$ nhân khẩu}
{$3700$ nhân khẩu}
\loigiai{
Tổng số dân: $1240 \times 3 = 3720$ người.
\textbf{Đáp án của câu hỏi là B}
}
\end{ex}

\begin{ex} %[3D2V56]
Tìm $x$, biết: $x : 3 = 1650$.
\choice
{$4850$}
{$5050$}
{\True $4950$}
{$4900$}
\loigiai{
$x = 1650 \times 3 = 4950$.
\textbf{Đáp án của câu hỏi là C}
}
\end{ex}

\begin{ex} %[3D2V56]
Một hình vuông có cạnh dài $1520\text{ m}$. Chu vi của hình vuông đó là
\choice
{\True $6080\text{ m}$}
{$6000\text{ m}$}
{$6120\text{ m}$}
{$5980\text{ m}$}
\loigiai{
Chu vi hình vuông: $1520 \times 4 = 6080\text{ (m)}$.
\textbf{Đáp án của câu hỏi là A}
}
\end{ex}

\begin{ex} %[3D2V56]
Một nhà máy mỗi ngày sản xuất được 1125 đôi giày. Hỏi trong 6 ngày làm việc nhà máy đó sản xuất được bao nhiêu đôi giày?
\choice
{$6650$ đôi giày}
{$6850$ đôi giày}
{$6700$ đôi giày}
{\True $6750$ đôi giày}
\loigiai{
Số đôi giày sản xuất được: $1125 \times 6 = 6750$ đôi.
\textbf{Đáp án của câu hỏi là D}
}
\end{ex}

\begin{ex} %[3D2V56]
Tính nhanh giá trị của biểu thức: $1250 \times 2 + 1250 \times 6$.
\choice
{$9000$}
{\True $10\,000$}
{$8000$}
{$12\,500$}
\loigiai{
$1250 \times (2 + 6) = 1250 \times 8 = 10\,000$.
\textbf{Đáp án của câu hỏi là B}
}
\end{ex}

\begin{ex} %[3D2V56]
Một xe tải chở 4 chuyến hàng, mỗi chuyến chở $2150\text{ kg}$ ngô. Người ta đã dỡ xuống $3600\text{ kg}$ ngô. Số ngô còn lại trên các xe là
\choice
{$4800\text{ kg}$}
{$5200\text{ kg}$}
{\True $5000\text{ kg}$}
{$4900\text{ kg}$}
\loigiai{
Tổng khối lượng ngô: $2150 \times 4 = 8600\text{ kg}$. Còn lại: $8600 - 3600 = 5000\text{ (kg)}$.
\textbf{Đáp án của câu hỏi là C}
}
\end{ex}

\begin{ex} %[3D2V56]
Tính giá trị của biểu thức: $A = 2314 \times 3 - 2800$.
\choice
{\True $4142$}
{$4132$}
{$4242$}
{$4042$}
\loigiai{
$2314 \times 3 = 6942$; $6942 - 2800 = 4142$.
\textbf{Đáp án của câu hỏi là A}
}
\end{ex}

\begin{ex} %[3D2V56]
Nếu một thừa số tăng thêm 2 đơn vị và thừa số kia giữ nguyên bằng 1450 thì tích tăng thêm bao nhiêu đơn vị?
\choice
{$1450$}
{$2800$}
{$3000$}
{\True $2900$}
\loigiai{
Tích tăng thêm: $1450 \times 2 = 2900$ đơn vị.
\textbf{Đáp án của câu hỏi là D}
}
\end{ex}

\begin{ex} %[3D2V56]
Người ta đóng 1200 gói kẹo vào các thùng, mỗi thùng 8 gói. Hỏi 7 thùng kẹo như thế có bao nhiêu gói kẹo? (Nếu mỗi thùng có 1200 gói kẹo thì 7 thùng có bao nhiêu gói kẹo?) Có 7 thùng kẹo, mỗi thùng chứa 1250 gói kẹo. Tổng số gói kẹo là
\choice
{$8650$ gói kẹo}
{\True $8750$ gói kẹo}
{$8850$ gói kẹo}
{$8500$ gói kẹo}
\loigiai{
Tổng số gói kẹo: $1250 \times 7 = 8750$ gói kẹo.
\textbf{Đáp án của câu hỏi là B}
}
\end{ex}

\begin{ex} %[3D2V56]
Tìm tích của số bé nhất có bốn chữ số khác nhau với 5.
\choice
{$5000$}
{$5120$}
{\True $5115$}
{$5125$}
\loigiai{
Số bé nhất có 4 chữ số khác nhau là 1023. Tích: $1023 \times 5 = 5115$.
\textbf{Đáp án của câu hỏi là C}
}
\end{ex}

\begin{ex} %[3D2V56]
Một cuộn dây dài $2450\text{ m}$. Bốn cuộn dây như thế dài tất cả bao nhiêu mét?
\choice
{\True $9800\text{ m}$}
{$9600\text{ m}$}
{$9900\text{ m}$}
{$9700\text{ m}$}
\loigiai{
$2450 \times 4 = 9800\text{ (m)}$.
\textbf{Đáp án của câu hỏi là A}
}
\end{ex}

\textbf{Phần 2. Trắc nghiệm đúng sai}
\setcounterpageref{ex}{0}

\begin{ex} %[3D2V56]
Xét tính đúng sai của phép nhân số có bốn chữ số với số có một chữ số:
\choiceTF
{\True Thực hiện phép nhân từ phải sang trái (bắt đầu từ hàng đơn vị).}
{$2500 \times 4 = 1000$.}
{\True Tích của số có 4 chữ số với 1 luôn bằng chính số đó.}
{Tích của $1500 \times 6$ bằng $8000$.}
\loigiai{
\begin{itemize}
    \item Ý a) Đúng: Quy tắc chuẩn.
    \item Ý b) Sai: Sai vì $2500 \times 4 = 10\,000$.
    \item Ý c) Đúng: Tính chất của 1.
    \item Ý d) Sai: Sai vì $1500 \times 6 = 9000$.
\end{itemize}
\begin{itemize}
    \item a) \textbf{Đúng}.
    \item b) \textbf{Sai}.
    \item c) \textbf{Đúng}.
    \item d) \textbf{Sai}.
\end{itemize}
}
\end{ex}

\begin{ex} %[3D2V56]
Một đội xe gồm 4 xe chở hàng, mỗi xe chở $1850\text{ kg}$ phân bón:
\choiceTF
{\True Đội xe chở tất cả $7400\text{ kg}$ phân bón.}
{\True Hai xe chở được $3700\text{ kg}$ phân bón.}
{Ba xe chở nhiều hơn $6000\text{ kg}$ phân bón.}
{\True Nếu mỗi bao nặng $50\text{ kg}$ thì mỗi xe chở 37 bao phân bón.}
\loigiai{
\begin{itemize}
    \item Ý a) Đúng: $1850 \times 4 = 7400\text{ kg}$.
    \item Ý b) Đúng: $1850 \times 2 = 3700\text{ kg}$.
    \item Ý c) Sai: Sai vì $1850 \times 3 = 5550\text{ kg} < 6000\text{ kg}$.
    \item Ý d) Đúng: $1850 : 50 = 37$ bao.
\end{itemize}
\begin{itemize}
    \item a) \textbf{Đúng}.
    \item b) \textbf{Đúng}.
    \item c) \textbf{Sai}.
    \item d) \textbf{Đúng}.
\end{itemize}
}
\end{ex}

\textbf{Phần 3. Trả lời ngắn}
\setcounterpageref{ex}{0}

\begin{ex} %[3D2V56]
Tính kết quả phép nhân: $1624 \times 5$.
\shortans{8120}
\loigiai{
\textbf{Đáp số:} $8120$\\
$1624 \times 5 = 8120$.
}
\end{ex}

\begin{ex} %[3D2V56]
Một hình vuông có cạnh dài $1850\text{ cm}$. Tính chu vi của hình vuông đó theo xăng-ti-mét.
\shortans{7400}
\loigiai{
\textbf{Đáp số:} $7400$\\
$1850 \times 4 = 7400\text{ (cm)}$.
}
\end{ex}

\begin{ex} %[3D2V56]
Tìm $x$, biết: $x : 4 = 1450$.
\shortans{5800}
\loigiai{
\textbf{Đáp số:} $5800$\\
$x = 1450 \times 4 = 5800$.
}
\end{ex}

\begin{ex} %[3D2V56]
Một xưởng may mỗi ngày may được 1250 chiếc khẩu trang. Hỏi trong 7 ngày xưởng đó may được bao nhiêu chiếc khẩu trang?
\shortans{8750}
\loigiai{
\textbf{Đáp số:} $8750$\\
$1250 \times 7 = 8750$ chiếc.
}
\end{ex}

\begin{ex} %[3D2V56]
Tính giá trị biểu thức: $1500 \times 6 - 3500$.
\shortans{5500}
\loigiai{
\textbf{Đáp số:} $5500$\\
$9000 - 3500 = 5500$.
}
\end{ex}

\begin{ex} %[3D2V56]
Có 3 bồn chứa nước, mỗi bồn chứa $2400\ell$ nước. Người ta đã dùng hết $4500\ell$. Hỏi trong 3 bồn còn lại tất cả bao nhiêu lít nước?
\shortans{2700}
\loigiai{
\textbf{Đáp số:} $2700$\\
Tổng: $2400 \times 3 = 7200\ell$. Còn lại: $7200 - 4500 = 2700\text{ (\ell)}$.
}
\end{ex}
